% chapters/09/section-01.tex
\section{Dari Dua Keadaan Menjadi Mesin Keadaan}
\label{sec:dua-keadaan}

NPC pada Bab~VIII hanya memiliki dua keadaan: mengejar dan berdiam. Bagian ini
menjelaskan apa yang ditambahkan ketika jumlahnya bertambah.

\subsection{Mesin Keadaan yang Paling Sederhana}

Perhatikan variabel \code{self.chasing} pada Bab~VIII. Variabel itu bernilai
benar atau salah, sehingga dengan sendirinya membagi seluruh perilaku NPC
menjadi dua golongan: perilaku ketika mengejar, dan perilaku ketika tidak
mengejar. Variabel itu sesungguhnya sudah merupakan mesin keadaan yang paling
sederhana, yaitu mesin dengan dua keadaan.

Yang belum ada hanyalah nama dan susunannya. Mesin keadaan yang sesungguhnya
memiliki tiga unsur yang tidak tampak pada variabel bernilai benar-salah.

\begin{enumerate}
  \item \textbf{Keadaan yang lebih dari dua}, masing-masing dengan perilaku
        yang berbeda. NPC yang berpatroli jelas berperilaku berbeda dari NPC
        yang berdiam, walaupun keduanya bukan sedang mengejar.
  \item \textbf{Syarat perpindahan yang tersendiri} untuk setiap pasangan
        keadaan. NPC berpindah dari berpatroli ke mengejar karena sebab yang
        berbeda dari sebab ia berpindah dari mengejar ke menyerang.
  \item \textbf{Pekerjaan penyiapan ketika sebuah keadaan dimasuki}, misalnya
        mengembalikan penghitung ke nol atau memilih titik patroli berikutnya.
\end{enumerate}

Bagian~2.5 telah menjelaskan gagasan mesin keadaan secara umum. Bagian ini dan
bagian berikutnya menerapkannya menjadi program.

\subsection{Empat Keadaan pada Bab Ini}

NPC pada bab ini memiliki empat keadaan, sesuai urutan yang disebut pada
RPS.

\begin{description}
  \item[berdiam] Berdiri tidak bergerak, sambil menghitung lama ia berdiam.
        Setelah 1,5 detik, ia mulai berpatroli.
  \item[berpatroli] Berpindah dari satu titik patroli ke titik berikutnya.
        Apabila pemain memasuki jari-jari persepsi, ia beralih mengejar.
  \item[mengejar] Bergerak menuju pemain dengan kecepatan tetap, sebagaimana
        dijelaskan pada Bab~VIII. Apabila pemain sudah cukup dekat, ia
        menyerang.
  \item[menyerang] Berhenti dan menyerang secara berkala. Apabila pemain
        menjauh melampaui jangkauan serang, ia kembali mengejar.
\end{description}

\subsection{Mengapa Tidak Cukup dengan Rangkaian Pemeriksaan}

Pertanyaan yang wajar muncul: mengapa tidak menuliskan seluruh perilaku itu
sebagai rangkaian pemeriksaan bersarang, tanpa memakai mesin keadaan sama
sekali?

\begin{kodekeluaran}
if jarak <= radius:
    if jarak <= jangkauan_serang:
        # menyerang
    else:
        # mengejar
else:
    if sudah_lama_berdiam:
        # berpatroli
    else:
        # berdiam
\end{kodekeluaran}

Cara itu memang berfungsi untuk empat keadaan. Akan tetapi, ia mengandung dua
kelemahan yang akan terasa ketika perilakunya bertambah.

\begin{itemize}
  \item \textbf{Pekerjaan penyiapan menjadi sulit ditempatkan.} Menentukan
        titik patroli berikutnya, atau mengembalikan pewaktu ke nol, harus
        ditulis di dalam cabang yang bersangkutan. Ketika jumlah keadaan
        bertambah, cabang-cabang itu menjadi panjang dan saling menyalin.
  \item \textbf{Keputusan dan tindakan tercampur.} Pada cara di atas, syarat
        perpindahan dan perilaku keadaan berada pada satu tempat. Akibatnya,
        syarat perpindahan tidak dapat diperiksa tanpa menjalankan perilakunya.
        Itulah yang membuat pengujian menjadi sulit, sebagaimana dibahas pada
        Bagian~\ref{sec:menguji-keadaan}.
\end{itemize}

Mesin keadaan memisahkan keduanya. Syarat perpindahan dikumpulkan pada satu
tempat, dan perilaku setiap keadaan dikumpulkan pada tempat lain. Pemisahan
itulah yang menjadi pokok dua bagian berikutnya.
